% --- Archivo: 03_convexidad.tex ---

% =========================================================
\section{Convexidad}
\label{v2:sec:1.3}
% =========================================================

\begin{definition}\label{v2:def:1.3.1}
    Un conjunto $D \subset \R^n$ es llamado \deftech{conjunto convexo} si para cualesquiera $x \in D$, $y \in D$ y $\alpha \in [0, 1]$, se tiene que $\alpha x + (1 - \alpha)y \in D$.
\end{definition}

Una \deftech{proyección} (ortogonal) del punto $x \in \R^n$ sobre un conjunto $D \subset \R^n$ es un punto de $D$ que está más próximo a $x$. En otras palabras, una proyección de $x$ sobre $D$ es una solución (global) del problema
\begin{equation}
    \min \|y - x\| \quad \text{sujeto a} \quad y \in D.
    \label{v2:eq:1.26}
\end{equation}

\begin{theorem}[Teorema de la proyección]\label{v2:thm:1.3.1}
    Si el conjunto $D$ es cerrado y convexo, entonces para todo $x \in \R^n$ la proyección de $x$ sobre $D$, denotada por $P_D(x)$, existe y es única. Además de eso, $\bar{x} = P_D(x)$ si, y solamente si,
    \begin{equation}
        \bar{x} \in D, \quad \langle x - \bar{x}, y - \bar{x} \rangle \le 0 \quad \forall y \in D.
        \label{v2:eq:1.27}
    \end{equation}
\end{theorem}

\begin{proof}
    Véase \cite{izmailov2020otimizacaovolume1} (o \cref{v1:cor:1.2.2} sobre proyección).
\end{proof}

\begin{definition}\label{v2:def:1.3.2}
    Sea $D \subset \R^n$ un conjunto convexo, se dice que la función $f : D \to \R$ es \deftech{convexa en $D$} cuando para cualesquiera $x \in D$, $y \in D$ y $\alpha \in [0, 1]$, se tiene
    \[
        f(\alpha x + (1 - \alpha)y) \le \alpha f(x) + (1 - \alpha)f(y).
    \]
\end{definition}

\begin{theorem}[Teorema de minimización convexa]\label{v2:thm:1.3.2}
    Sean $D \subset \R^n$ un conjunto convexo y $f : D \to \R$ una función convexa en $D$.
    Entonces todo minimizador local del problema es global. Además de eso, el conjunto de minimizadores es convexo.
\end{theorem}

Una función convexa es continua en cualquier subconjunto abierto de su dominio. Más aún, es localmente Lipschitz-continua en el interior de su dominio.

\begin{theorem}[Continuidad de funciones convexas]\label{v2:thm:1.3.3}
    Sean $\Omega \subset \R^n$ un conjunto convexo y abierto y $f : \Omega \to \R$ una función convexa en $\Omega$. Entonces $f$ es localmente Lipschitz-continua en $\Omega$.
\end{theorem}

\begin{proof}
    Véase \cite{izmailov2020otimizacaovolume1} (o \cref{v1:thm:3.4.2}).
\end{proof}

\begin{theorem}[Compacidad de conjuntos de nivel]\label{v2:thm:1.3.4}
    Supongamos que la función $f : \R^n \to \R$ sea fuertemente convexa en $\R^n$. Entonces el conjunto de nivel es compacto para todo $c \in \R$.
\end{theorem}

\begin{proof}
    Véase \cite{izmailov2020otimizacaovolume1} (o \cref{v1:thm:3.4.3}).
\end{proof}

\begin{corollary}\label{v2:cor:1.3.1}
    Sea $f : \R^n \to \R$ una función fuertemente convexa y $\emptyset \neq D \subset \R^n$ un conjunto cerrado cualquiera. Entonces $f$ tiene un minimizador en $D$ y él es único.
\end{corollary}

\begin{proof}
    Véase \cite{izmailov2020otimizacaovolume1} (o \cref{v1:cor:3.4.1}).
\end{proof}

Cuando una función es diferenciable, la convexidad admite varias caracterizaciones algebraicas.

\begin{theorem}[Caracterizaciones de funciones convexas diferenciables]\label{v2:thm:1.3.5}
    Sean $\Omega \subset \R^n$ un conjunto convexo y abierto y $f : \Omega \to \R$ una función diferenciable en $\Omega$. Entre otras cosas, $f$ es convexa si y solo si
    \[
        f(y) \ge f(x) + \langle f'(x), y - x \rangle.
    \]
\end{theorem}

\begin{proof}
    Véase \cite{izmailov2020otimizacaovolume1} (o \cref{v1:thm:3.4.5}).
\end{proof}

\begin{theorem}[Condiciones de optimalidad en caso convexo]\label{v2:thm:1.3.7}
    Sean $D \subset \R^n$ un conjunto convexo y cerrado, y $f : \R^n \to \R$ una función diferenciable en el punto $\bar{x} \in D$. Si $f$ es una función convexa y $\bar{x} = P_D(\bar{x} - \alpha f'(\bar{x}))$ para $\alpha > 0$, entonces $\bar{x}$ es un minimizador de $f$ en $D$.
\end{theorem}

\begin{proof}
    Véase \cite{izmailov2020otimizacaovolume1} (o \cref{v1:thm:3.4.6}).
\end{proof}

\begin{theorem}\label{v2:thm:1.3.8}
    Sea $\bar{x} \in D$ un punto estacionario del problema \eqref{v2:eq:1.14}, en el sentido de que valen las condiciones KKT. Sean $f$ y $g$ funciones convexas y sea $h$ una función afín. Entonces $\bar{x}$ es un minimizador del problema \eqref{v2:eq:1.14}.
\end{theorem}

\begin{proof}
    Véase \cite{izmailov2020otimizacaovolume1} (o \cref{v1:thm:4.2.2}).
\end{proof}

Recordamos ahora la noción de subdiferencial de una función convexa (no necesariamente diferenciable).

\begin{definition}\label{v2:def:1.3.4}
    Sea $f : \R^n \to \R$ una función convexa. Decimos que $y \in \R^n$ es un subgradiente de $f$ en el punto $x \in \R^n$ si
    \[
        f(z) \ge f(x) + \langle y, z - x \rangle \quad \forall z \in \R^n.
    \]
    El conjunto de todos los subgradientes de $f$ en $x$ se llama subdiferencial de $f$ en $x$; lo denotamos por $\partial f(x)$.
\end{definition}

\begin{theorem}[El subdiferencial de una función convexa]\label{v2:thm:1.3.9}
    Sea $f : \R^n \to \R$ una función convexa. Entonces para todo $x \in \R^n$, el conjunto $\partial f(x)$ es convexo, compacto y no vacío. Además, $f'(x; d) = \max \{ \langle y, d \rangle \mid y \in \partial f(x) \}$. El punto $\bar{x}$ es un minimizador de $f$ en $\R^n$ si, y solamente si, $0 \in \partial f(\bar{x})$.
\end{theorem}

\begin{proof}
    Véanse \cite{izmailov2020otimizacaovolume1} (o \cref{v1:thm:3.4.8,v1:thm:3.4.9}).
\end{proof}

\begin{definition}\label{v2:def:1.3.5}
    Sean $f : \R^n \to \R$ una función convexa y $\varepsilon \in \R_+$. Decimos que $y \in \R^n$ es un $\varepsilon$-subgradiente de $f$ en el punto $x \in \R^n$ cuando
    \[
        f(z) \ge f(x) + \langle y, z - x \rangle - \varepsilon \quad \forall z \in \R^n.
    \]
    El conjunto de todos los $\varepsilon$-subgradientes de $f$ en $x$, denotado por $\partial_\varepsilon f(x)$, se llama $\varepsilon$-subdiferencial de $f$ en $x$.
\end{definition}

\begin{theorem}[El $\varepsilon$-subdiferencial]\label{v2:thm:1.3.10}
    Sea $f : \R^n \to \R$ una función convexa. Entonces para todo $x \in \R^n$ y todo $\varepsilon \in \R_+$, el conjunto $\partial_\varepsilon f(x)$ es convexo, compacto y no vacío.
\end{theorem}

\begin{proof}
    Véase \cite{izmailov2020otimizacaovolume1} (o \cref{v1:def:3.4.2}).
\end{proof}